\documentclass[../../main.tex]{subfiles}
\begin{document}
\section{群作用与等变映射}\label{section:ch07-群作用与等变映射}

Lie 群对光滑流形理论最重要的应用涉及 Lie 群在其他流形上的作用。若 $G$ 是群，$M$ 是集合，则 $G$ 在 $M$ 上的\textbf{左作用}是映射 $G\times M\to M$，通常写作 $(g,p)\mapsto g\cdot p$，满足
\begin{align}
g_1\cdot(g_2\cdot p)=(g_1g_2)\cdot p\quad\text{对所有 }g_1,g_2\in G,\,p\in M;\qquad e\cdot p=p\quad\text{对所有 }p\in M.\label{align:ch07-06}
\end{align}
\textbf{右作用}类似地定义为映射 $M\times G\to M$，具有相应的复合法则：
\begin{align*}
(p\cdot g_1)\cdot g_2=p\cdot(g_1g_2)\quad\text{对所有 }g_1,g_2\in G,\,p\in M;\qquad p\cdot e=p\quad\text{对所有 }p\in M.
\end{align*}
若 $M$ 是拓扑空间，$G$ 是拓扑群，则 $G$ 在 $M$ 上的作用称为\textbf{连续作用}，若定义映射 $G\times M\to M$ 或 $M\times G\to M$ 是连续的。此时 $M$ 称为（左或右）$G$-空间。此外，若 $M$ 是带边或不带边的光滑流形，$G$ 是 Lie 群，且定义映射是光滑的，则该作用称为\textbf{光滑作用}。我们主要关注 Lie 群在光滑流形上的光滑作用。

有时给一个作用命名会很方便，例如 $\rho\colon G\times M\to M$；群元素 $g$ 对点 $p$ 的作用通常写作 $\rho_g(p)$。在这一记号下，左作用的条件\eqref{align:ch07-06}写作
\begin{align}
\rho_{g_1}\circ\rho_{g_2}=\rho_{g_1g_2},\qquad \rho_e=\operatorname{Id}_M;\label{align:ch07-07}
\end{align}
而对于右作用，第一个等式替换为
\begin{align*}
\rho_{g_2}\circ\rho_{g_1}=\rho_{g_1g_2}.
\end{align*}
对于光滑作用，每个映射 $\rho_g\colon M\to M$ 都是微分同胚，因为 $\rho_{g^{-1}}$ 是其光滑逆。

对于左作用，我们通常交替使用记号 $g\cdot p$ 和 $\rho_g(p)$。后一种记号包含更多信息，在需要明确指定所讨论的特定作用时很有用；而当前者所指作用已明确时，前者通常更方便。对于右作用，由于复合方式的原因，通常偏好记号 $p\cdot g$。

右作用总可以通过定义 $g\cdot p=p\cdot g^{-1}$ 的技巧转化为左作用，左作用也可以类似地转化为右作用。因此，关于左作用的任何结果都可以转化为关于右作用的结果，反之亦然。我们通常将注意力集中在左作用上，因为其群法则\eqref{align:ch07-07}具有群元素乘法对应于映射复合的性质。然而，在某些情形中右作用自然出现；我们将在本书后面看到几个这样的作用。

Lie 群作用通常出现在涉及某种对称性的情形中。例如，若 $M$ 是赋予度量或其他几何结构的向量空间或光滑流形，则 $M$ 上保持该结构的微分同胚的集合（称为该结构的\textbf{对称群}）往往是一个在 $M$ 上光滑作用的 Lie 群。

在全书中，我们使用如下关于群作用的标准术语。设 $\rho\colon G\times M\to M$ 是群 $G$ 在集合 $M$ 上的左作用。（右作用的定义类似。对于这些定义，不需要连续性或光滑性假设。）
\begin{itemize}
\item 对每个 $p\in M$，$p$ 的\textbf{轨道}，记作 $G\cdot p$，是 $p$ 在 $G$ 中元素作用下所有像的集合：
\begin{align*}
G\cdot p=\{g\cdot p:g\in G\}.
\end{align*}
\item 对每个 $p\in M$，$p$ 的\textbf{迷向群}或\textbf{稳定化子}，记作 $G_p$，是 $G$ 中固定 $p$ 的元素的集合：
\begin{align*}
G_p=\{g\in G:g\cdot p=p\}.
\end{align*}
群作用的定义保证 $G_p$ 是 $G$ 的子群。
\item 该作用称为\textbf{传递的}，若对每对点 $p,q\in M$，存在 $g\in G$ 使得 $g\cdot p=q$；等价地，若唯一的轨道就是整个 $M$。
\item 该作用称为\textbf{自由的}，若 $G$ 中固定 $M$ 中任意元素的唯一元素是单位元：对某个 $p\in M$，$g\cdot p=p$ 蕴涵 $g=e$；等价地，若每个迷向群都是平凡的。
\end{itemize}

下面是 Lie 群在流形上作用的一些例子。我们将在第 21 章中看到更多。

\begin{example}\label{example:ch07-22}[Lie 群作用]
\begin{enumerate}[(a)]
\item 若 $G$ 是任意 Lie 群，$M$ 是任意光滑流形，则 $G$ 在 $M$ 上的\textbf{平凡作用}定义为 $g\cdot p=p$，对所有 $g\in G$ 和 $p\in M$。它是光滑作用，每个轨道是单点，每个迷向群是整个 $G$。

\item $\operatorname{GL}(n,\mathbb{R})$ 在 $\mathbb{R}^n$ 上的\textbf{自然作用}是由矩阵乘法给出的左作用：$(A,x)\mapsto Ax$，将 $x\in\mathbb{R}^n$ 视为列矩阵。这是一个作用，因为 $I_n x=x$ 且矩阵乘法是结合的：$(AB)x=A(Bx)$。它是光滑的，因为 $Ax$ 的分量多项式地依赖于 $A$ 的矩阵元和 $x$ 的分量。由于任意非零向量可以被某个可逆线性变换映射到任意其他非零向量，恰有两个轨道：$\{0\}$ 和 $\mathbb{R}^n\setminus\{0\}$。

\item 每个 Lie 群 $G$ 通过左平移在自身上光滑作用。给定任意两点 $g_1,g_2\in G$，存在唯一的 $G$ 的左平移将 $g_1$ 映到 $g_2$，即由 $g_2g_1^{-1}$ 的左平移；因此该作用既是自由的又是传递的。更一般地，若 $H$ 是 $G$ 的 Lie 子群，则乘法映射在 $H\times G\to G$ 上的限制定义了 $H$ 在 $G$ 上的光滑且自由（但通常不传递）的左作用。类似的观察适用于右平移。

\item 每个 Lie 群通过共轭在自身上光滑作用：$g\cdot h=ghg^{-1}$。

\item 离散群 $\Gamma$ 在流形 $M$ 上的作用是光滑的，当且仅当对每个 $g\in\Gamma$，映射 $p\mapsto g\cdot p$ 是从 $M$ 到自身的光滑映射。因此，例如，$\mathbb{Z}^n$ 通过左平移在 $\mathbb{R}^n$ 上光滑且自由地作用：
\begin{align*}
(m^1,\dots,m^n)\cdot(x^1,\dots,x^n)=(m^1+x^1,\dots,m^n+x^n).
\end{align*}
\end{enumerate}
\end{example}

另一类重要的 Lie 群作用来自覆叠映射。设 $E$ 和 $M$ 是拓扑空间，$\pi\colon E\to M$ 是（拓扑）覆叠映射。$\pi$ 的一个\textbf{自同构}（也称为\textbf{覆盖变换}或\textbf{deck 变换}）是同胚 $\varphi\colon E\to E$ 使得 $\pi\circ\varphi=\pi$。所有自同构的集合 $\operatorname{Aut}_\pi(E)$ 称为 $\pi$ 的\textbf{自同构群}，在复合下构成群，从左侧作用于 $E$。可以证明，$\operatorname{Aut}_\pi(E)$ 在 $\pi$ 的每个纤维上传递作用当且仅当 $\pi$ 是正规覆叠映射，即对每个 $q\in E$，$\pi_*(\pi_1(E,q))$ 是 $\pi_1(M,\pi(q))$ 的正规子群（例如见 [LeeTM, Cor. 12.5]）。

\begin{proposition}\label{proposition:ch07-23}
设 $E$ 和 $M$ 是带边或不带边的光滑流形，$\pi\colon E\to M$ 是光滑覆叠映射。在离散拓扑下，自同构群 $\operatorname{Aut}_\pi(E)$ 是零维 Lie 群，在 $E$ 上光滑且自由地作用。
\end{proposition}

\begin{proof}
设 $\varphi\in\operatorname{Aut}_\pi(E)$ 是固定 $E$ 中点 $p$ 的自同构。只需旋转图表，就可以将 $\varphi$ 视为 $\pi$ 的提升。由于 $E$ 的恒等映射是另一个在 $p$ 处与 $\varphi$ 一致的提升，覆叠映射的唯一提升性质（\cref{proposition:appA-77}(a)）保证 $\varphi=\operatorname{Id}_E$。因此，$\operatorname{Aut}_\pi(E)$ 的作用是自由的。

为证明 $\operatorname{Aut}_\pi(E)$ 是 Lie 群，只需验证它是可数的。取任意 $q\in E$，令 $p=\pi(q)\in M$，$U\subseteq M$ 是 $p$ 的均匀覆叠邻域。由于 $E$ 是第二可数的，$\pi^{-1}(U)$ 有可数多个连通分支，且每个分支恰包含 $\pi^{-1}(p)$ 的一个点，因此 $\pi^{-1}(p)$ 是可数的。令 $\Phi_{(q)}\colon\operatorname{Aut}_\pi(E)\to E$ 为映射 $\Phi_{(q)}(\varphi)=\varphi(q)$。则 $\Phi_{(q)}$ 将 $\operatorname{Aut}_\pi(E)$ 映入 $\pi^{-1}(p)$，且作用是自由的这一事实蕴涵它是单射；因此 $\operatorname{Aut}_\pi(E)$ 是可数的。

作用的光滑性由\cref{theorem:ch04-29}可得。
\end{proof}

\textbf{等变映射}

对于某些具有群作用的流形，存在一个容易验证的充分条件，使得光滑映射具有常秩。设 $G$ 是 Lie 群，$M$ 和 $N$ 都是赋予光滑左或右 $G$-作用的光滑流形。映射 $F\colon M\to N$ 称为关于给定 $G$-作用\textbf{等变的}，若对每个 $g\in G$，
\begin{align*}
F(g\cdot p)=g\cdot F(p)\quad\text{（左作用）};\qquad F(p\cdot g)=F(p)\cdot g\quad\text{（右作用）}.
\end{align*}
等价地，若 $\rho$ 和 $\varphi$ 分别是 $M$ 和 $N$ 上给定的作用，则 $F$ 是等变的，当且仅当对每个 $g\in G$ 有如下图表交换：
\begin{align*}
\begin{tikzcd}[ampersand replacement=\&]
M \ar{r}{F} \ar{d}{\rho_g} \& N \ar{d}{\varphi_g}\\
M \ar{r}{F} \& N
\end{tikzcd}
\end{align*}
这个条件也表述为 $F$ \textbf{缠绕}$\rho$ 和 $\varphi$。

\begin{example}\label{example:ch07-24}
设 $v=(v^1,\dots,v^n)\in\mathbb{R}^n$ 是任意固定的非零向量。定义 $\mathbb{R}$ 在 $\mathbb{R}^n$ 和 $\mathbb{T}^n$ 上的光滑左作用为
\begin{align*}
t\cdot(x^1,\dots,x^n)=(x^1+tv^1,\dots,x^n+tv^n),\qquad
t\cdot(z^1,\dots,z^n)=\bigl(e^{2\pi\mathrm{i}tv^1}z^1,\dots,e^{2\pi\mathrm{i}tv^n}z^n\bigr),
\end{align*}
对 $t\in\mathbb{R}$，$(x^1,\dots,x^n)\in\mathbb{R}^n$，$(z^1,\dots,z^n)\in\mathbb{T}^n$。由 $\varepsilon^n(x^1,\dots,x^n)=(e^{2\pi\mathrm{i}x^1},\dots,e^{2\pi\mathrm{i}x^n})$ 给出的光滑映射 $\varepsilon^n\colon\mathbb{R}^n\to\mathbb{T}^n$ 关于这些作用是等变的。
\end{example}

下面\cref{theorem:ch07-05}的推广是证明某些映射具有常秩的极其有用的工具。

\begin{theorem}\label{theorem:ch07-25}[等变秩定理]
设 $M$ 和 $N$ 是光滑流形，$G$ 是 Lie 群。假设 $F\colon M\to N$ 是光滑映射，它关于 $M$ 上传递的光滑 $G$-作用和 $N$ 上任意光滑 $G$-作用是等变的。则 $F$ 具有常秩。因此，若 $F$ 是满射，则它是光滑淹没；若它是单射，则它是光滑浸入；若它是双射，则它是微分同胚。
\end{theorem}

\begin{figure}[ht]
\centering
\includegraphics[width=0.8\textwidth]{fig-07-3.png}
\caption{等变秩定理}\label{figure:ch07-equivariant-rank}
\end{figure}

\begin{proof}
记 $\rho$ 和 $\varphi$ 分别为 $M$ 和 $N$ 上的 $G$-作用，取 $M$ 中任意两点 $p$ 和 $q$。选取 $g\in G$ 使得 $\rho_g(p)=q$。（这样的 $g$ 存在，因为我们假设 $G$ 在 $M$ 上传递作用。）由于 $\varphi_g\circ F=F\circ\rho_g$，如下图表交换：
\begin{align*}
\begin{tikzcd}[ampersand replacement=\&]
T_pM \ar{r}{\mathrm{d}F_p} \ar{d}{\mathrm{d}(\rho_g)_p} \& T_{F(p)}N \ar{d}{\mathrm{d}(\varphi_g)_{F(p)}}\\
T_qM \ar{r}{\mathrm{d}F_q} \& T_{F(q)}N
\end{tikzcd}
\end{align*}
由于该图表中的纵向线性映射是同构，横向映射具有相同的秩。换言之，$F$ 在 $M$ 中任意两点 $p,q$ 处的秩相同，故 $F$ 具有常秩。最后一个结论由整体秩定理可得。
\end{proof}

下面是等变秩定理的一个重要应用。设 $G$ 是 Lie 群，$M$ 是光滑流形，$\rho\colon G\times M\to M$ 是光滑左作用。（右作用的定义类似。）对每个 $p\in M$，定义映射
\begin{align}
\rho^{(p)}\colon G\to M,\qquad \rho^{(p)}(g)=g\cdot p.\label{align:ch07-09}
\end{align}
这通常称为\textbf{轨道映射}，因为其像为轨道 $G\cdot p$。此外，原像 $\bigl(\rho^{(p)}\bigr)^{-1}(p)$ 是迷向群 $G_p$。

\begin{proposition}\label{proposition:ch07-26}[轨道映射的性质]
设 $\rho$ 是 Lie 群 $G$ 在光滑流形 $M$ 上的光滑左作用。对每个 $p\in M$，轨道映射 $\rho^{(p)}\colon G\to M$ 是光滑的且具有常秩，因此迷向群 $G_p=\bigl(\rho^{(p)}\bigr)^{-1}(p)$ 是 $G$ 的真嵌入 Lie 子群。若 $G_p=\{e\}$，则 $\rho^{(p)}$ 是单射光滑浸入，故轨道 $G\cdot p$ 是 $M$ 的浸入子流形。
\end{proposition}

\begin{remark}
事实上，每个轨道都是 $M$ 的浸入子流形，而不仅是对应于平凡迷向群的那些；但该事实的证明须等到第 21 章（见\cref{exercise:ch21-17}）。
\end{remark}

\begin{proof}
轨道映射是光滑的，因为它等于如下映射的复合：
\begin{align*}
G\xrightarrow{\iota}G\times\{p\}\hookrightarrow G\times M\xrightarrow{\rho}M.
\end{align*}
由群作用的定义可知，$\rho^{(p)}$ 关于 $G$ 通过左平移在自身上的作用和 $M$ 上给定的作用是等变的：
\begin{align*}
\rho^{(p)}(g_0g)=(g_0g)\cdot p=g_0\cdot(g\cdot p)=g_0\cdot\rho^{(p)}(g).
\end{align*}
由于 $G$ 在自身上传递作用，等变秩定理表明 $\rho^{(p)}$ 具有常秩。因此，由\cref{theorem:ch05-12}可知 $G_p$ 是真嵌入子流形，由\cref{proposition:ch07-11}可知是 Lie 子群。

现在假设 $G_p=\{e\}$。若 $\rho^{(p)}(g_0)=\rho^{(p)}(g)$，则
\begin{align*}
g_0\cdot p=g\cdot p\Rightarrow(g^{-1}g_0)\cdot p=p\Rightarrow g^{-1}g_0=e\Rightarrow g=g_0,
\end{align*}
表明 $\rho^{(p)}$ 是单射。由等变秩定理，它是光滑浸入，因此由\cref{proposition:ch05-18}可知轨道（赋予适当的拓扑和光滑结构）是浸入子流形。
\end{proof}

接下来我们利用等变秩定理来确定一般线性群的一些重要 Lie 子群。

\begin{example}\label{example:ch07-27}[正交群]
实 $n\times n$ 矩阵 $A$ 称为\textbf{正交的}，若作为线性映射 $A\colon\mathbb{R}^n\to\mathbb{R}^n$ 它保持 Euclid 点积：
\begin{align*}
(Ax)\cdot(Ay)=x\cdot y\quad\text{对所有 }x,y\in\mathbb{R}^n.
\end{align*}
所有正交 $n\times n$ 矩阵的集合 $O(n)$ 是 $\operatorname{GL}(n,\mathbb{R})$ 的子群，称为 $n$ 次\textbf{正交群}。容易验证，矩阵 $A$ 是正交的当且仅当它将 $\mathbb{R}^n$ 的标准基映为一组标准正交基，这等价于 $A$ 的列是标准正交的。由于矩阵 $A^{T}A$ 的 $(i,j)$ 元（其中 $A^{T}$ 表示 $A$ 的转置）是 $A$ 的第 $i$ 列和第 $j$ 列的点积，这个条件也等价于要求 $A^{T}A=I_n$。

定义光滑映射 $\Phi\colon\operatorname{GL}(n,\mathbb{R})\to\mathcal{M}(n,\mathbb{R})$ 为 $\Phi(A)=A^{T}A$。则 $O(n)$ 等于水平集 $\Phi^{-1}(I_n)$。为证明 $\Phi$ 具有常秩，从而 $O(n)$ 是嵌入 Lie 子群，我们证明 $\Phi$ 关于 $\operatorname{GL}(n,\mathbb{R})$ 的适当右作用是等变的。令 $\operatorname{GL}(n,\mathbb{R})$ 通过右乘法在自身上作用，并定义 $\operatorname{GL}(n,\mathbb{R})$ 在 $\mathcal{M}(n,\mathbb{R})$ 上的右作用为
\begin{align*}
X\cdot B=B^{T}XB\quad\text{对 }X\in\mathcal{M}(n,\mathbb{R}),\,B\in\operatorname{GL}(n,\mathbb{R}).
\end{align*}
容易验证这是光滑作用，且 $\Phi$ 是等变的，因为
\begin{align*}
\Phi(AB)=(AB)^{T}(AB)=B^{T}A^{T}AB=B^{T}\Phi(A)B=\Phi(A)\cdot B.
\end{align*}
因此，$O(n)$ 是 $\operatorname{GL}(n,\mathbb{R})$ 的真嵌入 Lie 子群。它是紧的，因为它在 $\mathcal{M}(n,\mathbb{R})\cong\mathbb{R}^{n^2}$ 中是闭且有界的：闭是因为它是 $\Phi$ 的水平集，有界是因为每个 $A\in O(n)$ 的列范数为 $1$，因此满足 $|A|=\sqrt{n}$。

为确定 $O(n)$ 的维数，需要计算 $\Phi$ 的秩。由于秩是常数，只需在单位矩阵 $I_n\in\operatorname{GL}(n,\mathbb{R})$ 处计算。因此对任意 $B\in T_{I_n}\operatorname{GL}(n,\mathbb{R})=\mathcal{M}(n,\mathbb{R})$，取曲线 $\gamma\colon(-\varepsilon,\varepsilon)\to\operatorname{GL}(n,\mathbb{R})$ 为 $\gamma(t)=I_n+tB$，计算
\begin{align*}
\mathrm{d}\Phi_{I_n}(B)=\left.\frac{\mathrm{d}}{\mathrm{d}t}\right|_{t=0}\Phi\circ\gamma(t)=\left.\frac{\mathrm{d}}{\mathrm{d}t}\right|_{t=0}(I_n+tB)^{T}(I_n+tB)=B^{T}+B.
\end{align*}
由这个公式，$\mathrm{d}\Phi_{I_n}$ 的像显然包含在对称矩阵的向量空间中。反过来，若 $B\in\mathcal{M}(n,\mathbb{R})$ 是任意对称 $n\times n$ 矩阵，则 $\mathrm{d}\Phi_{I_n}\bigl(\frac{1}{2}B\bigr)=B$。由此可知 $\mathrm{d}\Phi_{I_n}$ 的像恰为对称矩阵空间。这是 $\mathcal{M}(n,\mathbb{R})$ 的一个线性子空间，维数为 $n(n+1)/2$，因为每个对称矩阵由其主对角线及其上方的值唯一确定。因此，$O(n)$ 是维数为 $n^2-n(n+1)/2=n(n-1)/2$ 的嵌入 Lie 子群。
\end{example}

\begin{example}\label{example:ch07-28}[特殊正交群]
$n$ 次\textbf{特殊正交群}定义为 $\operatorname{SO}(n)=O(n)\cap\operatorname{SL}(n,\mathbb{R})\subseteq\operatorname{GL}(n,\mathbb{R})$。由于每个矩阵 $A\in O(n)$ 满足
\begin{align*}
1=\det I_n=\det(A^{T}A)=(\det A)(\det A^{T})=(\det A)^2,
\end{align*}
由此可知对所有 $A\in O(n)$ 有 $\det A=\pm1$。因此，$\operatorname{SO}(n)$ 是 $O(n)$ 中由行列式为正的矩阵组成的开子群，因而也是 $\operatorname{GL}(n,\mathbb{R})$ 中维数为 $n(n-1)/2$ 的嵌入 Lie 子群。它是紧群，因为它是 $O(n)$ 的闭子集。
\end{example}

\begin{example}\label{example:ch07-29}[酉群]
对任意复矩阵 $A$，$A$ 的\textbf{伴随}是矩阵 $A^*$，由 $A$ 的矩阵元取复共轭再转置构成：$A^*=\overline{A}^{T}$。注意 $(AB)^*=(\overline{A}\overline{B})^{T}=\overline{B}^{T}\overline{A}^{T}=B^{*}A^{*}$。对任意正整数 $n$，$n$ 次\textbf{酉群}是子群 $U(n)\subseteq\operatorname{GL}(n,\mathbb{C})$，由其列关于 Hermite 点积 $z\cdot w=\sum_i z_i\overline{w_i}$ 构成 $\mathbb{C}^n$ 的标准正交基的复 $n\times n$ 矩阵组成。容易验证 $U(n)$ 由满足 $A^{*}A=I_n$ 的矩阵 $A$ 组成。\cref{exercise:ch07-13}表明它是 $\operatorname{GL}(n,\mathbb{C})$ 的维数为 $n^2$ 的真嵌入 Lie 子群。
\end{example}

\begin{example}\label{example:ch07-30}[特殊酉群]
群 $\operatorname{SU}(n)=U(n)\cap\operatorname{SL}(n,\mathbb{C})$ 称为 $n$ 次\textbf{复特殊酉群}。\cref{exercise:ch07-14}表明它是 $U(n)$ 的维数为 $n^2-1$ 的真嵌入 Lie 子群。由于光滑嵌入的复合 $\operatorname{SU}(n)\hookrightarrow U(n)\hookrightarrow\operatorname{GL}(n,\mathbb{C})$ 仍是光滑嵌入，这蕴涵 $\operatorname{SU}(n)$ 也嵌入在 $\operatorname{GL}(n,\mathbb{C})$ 中。
\end{example}

\textbf{半直积}

群作用为我们提供了构造 Lie 群的强有力的新方法。设 $H$ 和 $N$ 是 Lie 群，$\rho\colon H\times N\to N$ 是 $H$ 在 $N$ 上的光滑左作用。若对每个 $h\in H$，映射 $\rho_h\colon N\to N$ 是 $N$ 的群自同构（即从 $N$ 到自身的同构），则称该作用是\textbf{通过自同构的作用}。给定这样的作用，我们定义一个新的 Lie 群 $N\rtimes_\rho H$，称为 $H$ 和 $N$ 的\textbf{半直积}，如下所示。作为光滑流形，$N\rtimes_\rho H$ 就是 Cartesian 积 $N\times H$；但群乘法定义为
\begin{align}
(n,h)(n',h')=\bigl(n\rho_h(n'),hh'\bigr).\label{align:ch07-10}
\end{align}
有时，若 $H$ 在 $N$ 上的作用是已知的或无关紧要的，半直积简记为 $N\rtimes H$。

\begin{exercise}\label{exercise:ch07-31}
验证\eqref{align:ch07-10}确实在流形 $N\times H$ 上定义了 Lie 群结构，以 $(e,e)$ 为单位元，且 $(n,h)^{-1}=\bigl(\rho_{h^{-1}}(n^{-1}),h^{-1}\bigr)$。
\end{exercise}

\begin{example}\label{example:ch07-32}[Euclid 群]
若将 $\mathbb{R}^n$ 视为加法下的 Lie 群，则 $O(n)$ 在 $\mathbb{R}^n$ 上的自然作用是通过自同构的作用。所得半直积 $E(n)=\mathbb{R}^n\rtimes O(n)$ 称为\textbf{Euclid 群}；其乘法由 $(b,A)(b',A')=(b+Ab',AA')$ 给出。它通过 $(b,A)\cdot x=b+Ax$ 在 $\mathbb{R}^n$ 上作用。该作用保持直线、距离和角度度量，因而保持 Euclid 几何的所有关系。
\end{example}

下一个命题详述半直积的一些基本性质。回忆\eqref{align:ch07-05}中为群 $G$ 的子集 $A,B$ 定义的记号 $AB$。

\begin{proposition}\label{proposition:ch07-33}[半直积的性质]
设 $N$ 和 $H$ 是 Lie 群，$\rho$ 是 $H$ 在 $N$ 上通过自同构的光滑作用。令 $G=N\rtimes_\rho H$。
\begin{enumerate}[(a)]
\item 子集 $\widetilde{N}=N\times\{e\}$ 和 $\widetilde{H}=\{e\}\times H$ 是 $G$ 的闭 Lie 子群，分别同构于 $N$ 和 $H$。
\item $\widetilde{N}$ 是 $\widetilde{G}$ 的正规子群。
\item $\widetilde{N}\cap\widetilde{H}=\{(e,e)\}$ 且 $\widetilde{N}\widetilde{H}=G$。
\end{enumerate}
\end{proposition}

\begin{exercise}\label{exercise:ch07-34}
证明上述命题。
\end{exercise}

由于下一个定理，许多 Lie 群可以实现为适当子群的半直积。

\begin{theorem}\label{theorem:ch07-35}[半直积的刻画]
设 $G$ 是 Lie 群，$N,H\subseteq G$ 是闭 Lie 子群，满足 $N$ 是正规的，$N\cap H=\{e\}$，且 $NH=G$。则映射 $(n,h)\mapsto nh$ 是 $N\rtimes_\rho H$ 与 $G$ 之间的 Lie 群同构，其中 $\rho\colon H\times N\to N$ 是通过共轭的作用：$\rho_h(n)=hnh^{-1}$。
\end{theorem}

\begin{proof}
见\cref{exercise:ch07-18}。
\end{proof}

在\cref{theorem:ch07-35}的假设下，我们称 $G$ 是 $N$ 和 $H$ 的\textbf{内半直积}。一些例子在\cref{exercise:ch07-20}中描述。

\textbf{表示}

到目前为止我们见过的大多数 Lie 群都可以实现为 $\operatorname{GL}(n,\mathbb{R})$ 或 $\operatorname{GL}(n,\mathbb{C})$ 的 Lie 子群。自然要问是否所有 Lie 群都具有这种形式。研究这个问题的关键是群表示理论。

回忆，若 $V$ 是有限维实或复向量空间，$\operatorname{GL}(V)$ 表示 $V$ 的可逆线性变换群，当 $n=\dim V$ 时它是同构于 $\operatorname{GL}(n,\mathbb{R})$ 或 $\operatorname{GL}(n,\mathbb{C})$ 的 Lie 群。若 $G$ 是 Lie 群，则 $G$ 的一个（有限维）\textbf{表示}是从 $G$ 到某个 $V$ 的 $\operatorname{GL}(V)$ 的 Lie 群同态。（虽然对许多应用而言考虑 $V$ 为无限维的情形也很有用，但在本书中我们只考虑有限维表示。）

若表示 $\rho\colon G\to\operatorname{GL}(V)$ 是单射的，则称为\textbf{忠实的}。在这种情况下，由\cref{proposition:ch07-17}可知 $\rho$ 的像是 $\operatorname{GL}(V)$ 的 Lie 子群，且 $\rho$ 给出 $G$ 与 $\rho(G)\subseteq\operatorname{GL}(V)\cong\operatorname{GL}(n,\mathbb{R})$ 或 $\operatorname{GL}(n,\mathbb{C})$ 之间的 Lie 群同构。因此，Lie 群容许忠实表示当且仅当它同构于某个 $\operatorname{GL}(n,\mathbb{R})$ 或 $\operatorname{GL}(n,\mathbb{C})$ 的 Lie 子群。并非每个 Lie 群都容许这样的表示。我们尚未具备构造反例的技术，但\cref{exercise:ch21-26}要求你证明 $\operatorname{SL}(2,\mathbb{R})$ 的万有覆叠群没有忠实表示，因而不同构于任何矩阵群。

表示论是一门广泛的学科，其应用涉及微分几何、微分方程、调和分析、数论、量子物理和工程等众多领域；我们在此只能略作涉及。

\begin{example}\label{example:ch07-36}[Lie 群表示]
\begin{enumerate}[(a)]
\item 若 $G$ 是 $\operatorname{GL}(n,\mathbb{R})$ 的任意 Lie 子群，则包含映射 $G\hookrightarrow\operatorname{GL}(n,\mathbb{R})=\operatorname{GL}(\mathbb{R}^n)$ 是忠实表示，称为 $G$ 的\textbf{定义表示}。$\operatorname{GL}(n,\mathbb{C})$ 的 Lie 子群的定义表示类似定义。

\item 包含映射 $S^1\hookrightarrow\mathbb{C}^*\cong\operatorname{GL}(1,\mathbb{C})$ 是圆周群的忠实表示。更一般地，映射 $\rho\colon\mathbb{T}^n\to\operatorname{GL}(n,\mathbb{C})$ 由
\begin{align*}
\rho(z^1,\dots,z^n)=\begin{pmatrix}
z^1&0&\cdots&0\\
0&z^2&\cdots&0\\
\vdots&\vdots&\ddots&\vdots\\
0&0&\cdots&z^n
\end{pmatrix}
\end{align*}
给出，是 $\mathbb{T}^n$ 的忠实表示。

\item 令 $\sigma\colon\mathbb{R}^n\to\operatorname{GL}(n+1,\mathbb{R})$ 为将 $x\in\mathbb{R}^n$ 映为矩阵 $\sigma(x)$ 的映射，以分块形式定义为
\begin{align*}
\sigma(x)=\begin{pmatrix}I_n&x\\ 0&1\end{pmatrix},
\end{align*}
其中 $I_n$ 是 $n\times n$ 单位矩阵，$x$ 视为 $n\times1$ 列矩阵。直接计算表明 $\sigma$ 是加法 Lie 群 $\mathbb{R}^n$ 的忠实表示。

\item $\mathbb{R}^n$ 的另一个忠实表示是映射 $\mathbb{R}^n\to\operatorname{GL}(n,\mathbb{R})$，将 $(x^1,\dots,x^n)$ 映为对角元为 $(e^{x^1},\dots,e^{x^n})$ 的对角矩阵。

\item $\mathbb{R}^n$ 的又一个表示是映射 $\mathbb{R}^n\to\operatorname{GL}(n,\mathbb{C})$，将 $x$ 映为对角元为 $(e^{2\pi\mathrm{i}x^1},\dots,e^{2\pi\mathrm{i}x^n})$ 的对角矩阵。这个表示不是忠实的，因为其核是子群 $\mathbb{Z}^n\subseteq\mathbb{R}^n$。

\item 设 $E(n)$ 为 Euclid 群（\cref{example:ch07-32}）。$E(n)$ 的一个忠实表示由映射 $\rho\colon E(n)\to\operatorname{GL}(n+1,\mathbb{R})$ 给出，以分块形式定义为
\begin{align*}
\rho(b,A)=\begin{pmatrix}A&b\\ 0&1\end{pmatrix},
\end{align*}
其中 $b$ 视为列矩阵。

\item 对正整数 $n$ 和 $d$，记 $\mathcal{P}_d^n$ 为次数至多为 $d$ 的实值多项式函数 $p\colon\mathbb{R}^n\to\mathbb{R}$ 的向量空间。对任意矩阵 $A\in\operatorname{GL}(n,\mathbb{R})$，定义线性映射 $\tau_d^n(A)\colon\mathcal{P}_d^n\to\mathcal{P}_d^n$ 为
\begin{align*}
\tau_d^n(A)p=p\circ A^{-1}.
\end{align*}
\cref{exercise:ch07-24}表明映射 $\tau_d^n\colon\operatorname{GL}(n,\mathbb{R})\to\operatorname{GL}(\mathcal{P}_d^n)$ 是忠实表示。
\end{enumerate}
\end{example}

表示与群作用之间有密切联系。设 $G$ 是 Lie 群，$V$ 是有限维向量空间。$G$ 在 $V$ 上的作用称为\textbf{线性作用}，若对每个 $g\in G$，由 $x\mapsto g\cdot x$ 给出的从 $V$ 到自身的映射是线性的。例如，若 $\rho\colon G\to\operatorname{GL}(V)$ 是 $G$ 的表示，则有 $G$ 在 $V$ 上相伴的光滑线性作用，由 $g\cdot x=\rho(g)x$ 给出。下一个命题表明每个线性作用都属于这种类型。

\begin{proposition}\label{proposition:ch07-37}
设 $G$ 是 Lie 群，$V$ 是有限维向量空间。$G$ 在 $V$ 上的光滑左作用是线性的，当且仅当它形如 $g\cdot x=\rho(g)x$，其中 $\rho$ 是 $G$ 的某个表示。
\end{proposition}

\begin{proof}
由表示诱导的每个作用显然是线性的。为证明逆命题，假设给定 $G$ 在 $V$ 上的线性作用。假设蕴涵对每个 $g\in G$，存在线性映射 $\rho(g)\in\operatorname{GL}(V)$ 使得对所有 $x\in V$ 有 $g\cdot x=\rho(g)x$。作用满足\eqref{align:ch07-06}这一事实转化为 $\rho(g_1g_2)=\rho(g_1)\rho(g_2)$，故 $\rho\colon G\to\operatorname{GL}(V)$ 是群同态。因此，为证明它是 Lie 群表示，只需证明它是光滑的。

选取 $V$ 的一组基 $(E_i)$，对每个 $i$，令 $\pi_i\colon V\to\mathbb{R}$ 为关于该基到第 $i$ 个坐标的投影：$\pi_i(\sum_j x^j E_j)=x^i$。若记 $\rho_i^j(g)$ 为 $\rho(g)$ 关于该基的矩阵元，则 $\rho_i^j(g)=\pi_i(g\cdot E_j)$，故每个函数 $\rho_i^j$ 是光滑函数的复合。由于矩阵元构成 $\operatorname{GL}(V)$ 的整体光滑坐标，这蕴涵 $\rho$ 是光滑的。
\end{proof}

\end{document}